% SEGMENT OLP-0061-S01
% Part: propositional-logic
% Chapter: syntax-and-semantics
% Section: valuations-sat

\documentclass[../../../include/open-logic-section]{subfiles}

\begin{document}

\olfileid{pl}{syn}{val}

\olsection{\usetoken{P}{valuation} மற்றும் நிறைவுறுத்தல்}

\begin{defn}[!!^{valuation}s]
``மெய்'', ``பொய்'' என்ற இரு மெய்மதிப்புகளின் கணம் $\{\True, \False\}$
என்க. $\Lang{L_0}$ க்கான \emph{!!{valuation}} என்பது மொழியின்
!!{propositional variable}s க்கு $\True$ அல்லது $\False$ ஐ ஒதுக்கும்
சார்பு~$\pAssign{v}$; அதாவது $\pAssign{v} \colon
\PVar \to \{\True, \False \}$.
\end{defn}

% SEGMENT OLP-0061-S02
\begin{defn}
  !!a{valuation}~$\pAssign{v}$ தரப்பட்டால், மதிப்பீட்டுச் சார்பு
  $\pValue{v} \colon \Frm[L_0] \to \{\True, \False \}$ ஐப் பின்வருமாறு
  தொகுத்தறிதலாக வரையறுக்க:
  \begin{align*}
    \iftag{prvFalse}{\pValue{v}(\lfalse) & = \False; \\}{}
    \iftag{prvTrue}{\pValue{v}(\ltrue) & = \True; \\}{}
    \pValue{v}(\Obj p_n) & = \pAssign{v}(\Obj p_n); \\
    \iftag{prvNot}{\pValue{v}(\lnot !A) & = \begin{cases}
        \True & \text{பின்வரும் நிலையில்: } \pValue{v}(!A) = \False;\\
        \False & \text{இல்லையெனில்.}
      \end{cases} \\ }{}
    \iftag{prvAnd}{\pValue{v}(!A \land !B) & = \begin{cases}
        \True &
        \text{$\pValue{v}(!A) = \True$ மற்றும் $\pValue{v}(!B) = \True$ எனில்;}\\
        \False &
        \text{$\pValue{v}(!A) = \False$ அல்லது $\pValue{v}(!B) = \False$ எனில்}.
      \end{cases}\\}{}
    \iftag{prvOr}{\pValue{v}(!A \lor !B) & = \begin{cases}
        \True &
        \text{$\pValue{v}(!A) = \True$ அல்லது $\pValue{v}(!B) = \True$ எனில்;}\\
        \False &
        \text{$\pValue{v}(!A) = \False$ மற்றும் $\pValue{v}(!B) = \False$ எனில்}.
      \end{cases}\\}{}
    \iftag{prvIf}{\pValue{v}(!A \lif !B) & = \begin{cases}
        \True &
        \text{$\pValue{v}(!A) = \False$ அல்லது $\pValue{v}(!B) = \True$ எனில்;}\\
        \False &
        \text{$\pValue{v}(!A) = \True$ மற்றும் $\pValue{v}(!B) = \False$ எனில்}.
      \end{cases}\\}{}
    \iftag{prvIff}{\pValue{v}(!A \liff !B) & = \begin{cases}
        \True &
        \text{$\pValue{v}(!A) = \pValue{v}(!B)$ எனில்;}\\
        \False &
        \text{$\pValue{v}(!A) \neq \pValue{v}(!B)$ எனில்}.
      \end{cases}}{}
\end{align*}
\end{defn}

% SEGMENT OLP-0061-S03
\begin{explain}
இந்தக் கூறுகள் பின்வரும் மெய்மை அட்டவணைகளுக்கு ஒத்தவை:
\begin{center}
\iftag{prvNot}{
  \begin{tabular}{|c||c|} \hline
    $!A$ & $ \lnot !A$ \\
    \hline \hline
    $\True$ & $\False$ \\
    $\False$ & $\True$ \\
    \hline
  \end{tabular}
}{}
\iftag{prvAnd}{
  \begin{tabular}{|cc||c|} \hline
    $!A$ & $!B$ & $!A \land !B$ \\
    \hline \hline
    $\True$ & $\True$ & $\True$ \\
    $\True$ & $\False$ & $\False$ \\
    $\False$ & $\True$ & $\False$ \\
    $\False$ & $\False$ & $\False$ \\
    \hline
  \end{tabular}
}{}
\iftag{prvOr}{
  \begin{tabular}{|cc||c|} \hline
    $!A$ & $!B$ & $!A \lor !B$ \\
    \hline \hline
    $\True$ & $\True$ & $\True$ \\
    $\True$ & $\False$ & $\True$ \\
    $\False$ & $\True$ & $\True$ \\
    $\False$ & $\False$ & $\False$ \\
    \hline
  \end{tabular}
}{}%
\iftag{notprvNot,notprvAnd,notprvOr,notprvIf}{}{\\[1em]}
\iftag{prvIf}{
  \begin{tabular}{|cc||c|} \hline
    $!A$ & $!B$ & $!A \lif !B$ \\
    \hline \hline
    $\True$ & $\True$ & $\True$ \\
    $\True$ & $\False$ & $\False$ \\
    $\False$ & $\True$ & $\True$ \\
    $\False$ & $\False$ & $\True$ \\
    \hline
  \end{tabular}
}{}
\iftag{prvIff}{
  \begin{tabular}{|cc||c|} \hline
    $!A$ & $!B$ & $!A \liff !B$ \\
    \hline \hline
    $\True$ & $\True$ & $\True$ \\
    $\True$ & $\False$ & $\False$ \\
    $\False$ & $\True$ & $\False$ \\
    $\False$ & $\False$ & $\True$ \\
    \hline
  \end{tabular}
}{}
\end{center}
\end{explain}

% SEGMENT OLP-0061-S04
\begin{prob}
$\Lang{L_0}$ இல் $\diamondsuit$ என்ற மூவிட இணைப்பைப் பின்வரும்
மதிப்பீட்டுடன் சேர்ப்பதைக் கருதுக:
\begin{gather*}
  \pValue{v}(\diamondsuit( !A , !B , !C ) ) = \begin{cases}
    \pValue{v}( !B ) &
    \text{$\pValue{v}(!A) = \True$ எனில்;}\\
    \pValue{v}( !C ) &
    \text{$\pValue{v}(!A) = \False $ எனில்.}
  \end{cases}
\end{gather*}
இந்த இணைப்பிற்கான மெய்மை அட்டவணையை எழுதுக.
\end{prob}

% SEGMENT OLP-0061-S05
\begin{thm}[உள்ளூர்த் தீர்மானிப்பு]
\ollabel{thm:LocalDetermination} !!{valuation}s $\pAssign{v_1}$,
$\pAssign{v_2}$ ஆகியவை $!A$ இல் வரும் !!{propositional
variable}s
மீது ஒத்திருக்கட்டும். அதாவது, ஏதோ ஒரு !!{formula}~$!A$ இல்
$\Obj p_n$ வரும்போதெல்லாம் $\pAssign{v_1}(\Obj p_n) =
\pAssign{v_2}(\Obj p_n)$ என இருக்கட்டும். அப்போது $\pValue{v_1}$,
$\pValue{v_2}$ ஆகியவையும் $!A$ மீது ஒத்திருக்கும்; அதாவது
$\pValue{v_1}(!A) = \pValue{v_2}(!A)$.
\end{thm}

\begin{proof}
$!A$ இன்மீதான தொகுத்தறிதலால் நிறுவுக.
\end{proof}

% SEGMENT OLP-0061-S06
\begin{defn}[நிறைவுறுத்தல்]
\ollabel{defn:satisfaction} \emph{!!a{formula}~$!A$ ஐ
  !!a{valuation}~$\pAssign{v}$ நிறைவுறுத்துதல்} என்ற கருத்தை
  $\pSat{v}{!A}$ எனப் பின்வருமாறு தொகுத்தறிதலாக வரையறுக்கலாம்.
  (``$\pSat{v}{!A}$ அன்று'' என்பதைக் குறிக்க $\pSat/{v}{!A}$ என்று
  எழுதுகிறோம்.)
\begin{enumerate}
\tagitem{prvFalse}{%
  \indcase{!A}{\lfalse}{$\pSat/{v}{\indfrm}$.}}{}

\tagitem{prvTrue}{%
  \indcase{!A}{\ltrue}{$\pSat{v}{\indfrm}$.}}{}

\item \indcase{!A}{\Obj p_i}{$\pSat{v}{\indfrm}$ என்பது, அப்போதும்
  அப்போது மட்டுமே, $\pAssign{v}(\Obj p_i) = \True$.}

\tagitem{prvNot}{%
  \indcase{!A}{\lnot !B}{$\pSat{v}{\indfrm}$ என்பது, அப்போதும் அப்போது
    மட்டுமே, $\pSat/{v}{!B}$.}}{}

\tagitem{prvAnd}{%
  \indcase{!A}{(!B \land !C)}{$\pSat{v}{\indfrm}$ என்பது, அப்போதும்
    அப்போது மட்டுமே, $\pSat{v}{!B}$ மற்றும் $\pSat{v}{!C}$.}}{}

\tagitem{prvOr}{%
  \indcase{!A}{(!B \lor !C)}{$\pSat{v}{\indfrm}$ என்பது, அப்போதும்
    அப்போது மட்டுமே, $\pSat{v}{!B}$ அல்லது $\pSat{v}{!C}$ (அல்லது
    இரண்டும்).}}{}

\tagitem{prvIf}{%
  \indcase{!A}{(!B \lif !C)}{$\pSat{v}{\indfrm}$ என்பது, அப்போதும்
    அப்போது மட்டுமே, $\pSat/{v}{!B}$ அல்லது $\pSat{v}{!C}$ (அல்லது
    இரண்டும்).}}{}

\tagitem{prvIff}{%
  \indcase{!A}{(!B \liff !C)}{$\pSat{v}{\indfrm}$ என்பது, அப்போதும்
    அப்போது மட்டுமே, $\pSat{v}{!B}$ மற்றும் $\pSat{v}{!C}$ இரண்டும்
    பொருந்தும், அல்லது $\pSat{v}{!B}$ மற்றும் $\pSat{v}{!C}$ இரண்டுமே
    பொருந்தாது.}}{}
\end{enumerate}
$\Gamma$ என்பது !!{formula}s இன் கணம் எனில், ஒவ்வொரு~$!A \in \Gamma$
க்கும் $\pSat{v}{!A}$ என இருப்பது, அப்போதும் அப்போது மட்டுமே,
$\pSat{v}{\Gamma}$ என்க.
\end{defn}

% SEGMENT OLP-0061-S07
\begin{prop}\ollabel{prop:sat-value}
  $\pSat{v}{!A}$ என்பது, அப்போதும் அப்போது மட்டுமே,
  $\pValue{v}(!A) = \True$.
\end{prop}

\begin{proof}
  $!A$ இன்மீதான தொகுத்தறிதலால் நிறுவுக.
\end{proof}

\begin{prob}
  \olref[pl][syn][val]{prop:sat-value} ஐ நிறுவுக.
\end{prob}

\end{document}
